
\documentclass[uplatex,a4paper,11pt]{jsarticle}
\usepackage[top=25mm,bottom=25mm,left=25mm,right=25mm]{geometry}
\usepackage{longtable,booktabs,array,fancyhdr,lastpage}
\usepackage[dvipdfmx]{hyperref}
\usepackage{xcolor}
\hypersetup{colorlinks=true,linkcolor=black,urlcolor=black}
\renewcommand{\baselinestretch}{1.18}
\pagestyle{fancy}
\fancyhf{}
\lhead{08 章 ソフトウェア構成管理}
\rhead{SWEBOK v4.0a 日本語まとめ}
\cfoot{\thepage/\pageref{LastPage}}
\newcommand{\noteBox}[2]{\par\noindent\fbox{\begin{minipage}{0.96\linewidth}\textbf{#1}\par\medskip #2\end{minipage}}\par\medskip}
\newcommand{\todobox}[2]{\par\noindent\fbox{\begin{minipage}{0.96\linewidth}\textbf{TODO（図） #1}\par\medskip\textbf{画像生成用プロンプト}\par #2\end{minipage}}\par\medskip}
\begin{document}
\begin{titlepage}
\thispagestyle{empty}
\vspace*{55mm}
\begin{center}
{\Large\bfseries 第 08 章 ソフトウェア構成管理}\par\vspace{12mm}
{Software Configuration Management の日本語まとめ}\par\vspace{9mm}
{SWEBOK Guide v4.0a Chapter 08 を初学者向けに再構成}\par\vspace{14mm}
{作成日：2026 年 5 月 12 日}\par
\end{center}
\vspace{16mm}\hrule\vspace{10mm}
\noindent\textbf{この資料の主張}\par
ソフトウェア構成管理は、単なる Git の使い方やファイル名の規則ではない。ソフトウェアを構成する品目を識別し、変更を統制し、状態を記録し、監査し、正しい内容をビルドしてリリースするための管理活動である。ソフトウェアは変更しやすいが、変更しやすいからこそ、何を、誰が、いつ、なぜ、どの版に対して変えたのかを追跡できなければ、品質、セキュリティ、保守性、説明責任が崩れる。
\noteBox{この資料の読み方}{専門用語は原則として日本語で示し、必要に応じて英語や略語を括弧内に併記する。英語名は、元資料やツール上の名称を確認したいときの手がかりとして残す。初学者が前提知識なしに読めるように、各節では「何を管理するのか」「なぜ管理するのか」「実務では何を確認するのか」を重視する。}
\end{titlepage}
\setcounter{page}{1}
\tableofcontents
\clearpage
\section*{全体概要：この章で学ぶこと}
\addcontentsline{toc}{section}{全体概要：この章で学ぶこと}
この章は、ソフトウェアを構成する品目を識別し、変更を統制し、状態を記録し、監査し、正しい内容をビルドしてリリースするための知識を扱う。ソフトウェア構成管理は、版管理ツールの操作だけではなく、ソフトウェアの完全性、正確性、追跡可能性、説明責任を保つための管理体系である。

\noteBox{到達目標}{この章を読み終えると、構成品目、基準版、変更要求、構成状態記録、構成監査、リリース管理の役割を説明できるようになる。また、なぜ構成管理が品質保証、セキュリティ、保守、規制対応に関わるのかを説明できるようになる。}
\begin{itemize}
\item 「何を管理対象にするか」を構成品目として定められる。
\item 構成品目を変更する前に、変更要求、影響分析、承認、実施、検証の流れを説明できる。
\item 基準版が、変更を止めるものではなく、変更を管理可能にする基点であることを説明できる。
\item 状態記録・報告が、証跡、品質保証、セキュリティ、規制対応に使われることを説明できる。
\item ビルド、部品表、リリースノート、版記述文書がリリース管理で果たす役割を説明できる。
\end{itemize}
\todobox{第8章の全体地図を入れる}{白背景、モノクロ、A4 本文幅に収まる横長の学習用図解を作成する。中央に「ソフトウェア構成管理」を置き、周囲に「プロセス管理」「構成識別」「変更統制」「状態記録・報告」「構成監査」「リリース管理と提供」「構成管理ツール」を配置する。矢印で「識別した品目を変更統制し、状態を記録し、監査し、リリースへつなげる」流れを示す。ラベルは日本語、細線、講義ノート風、余白を広めにする。}
\section*{最初に必要な用語}
\addcontentsline{toc}{section}{最初に必要な用語}
構成管理では略語が多い。初学者は、まず「構成品目」「基準版」「変更要求」「状態記録」「監査」の五つを軸に読むとよい。

\begin{longtable}{p{0.30\linewidth}p{0.64\linewidth}}\toprule
用語 & 初学者向けの説明 \\ \midrule
\endfirsthead\toprule
用語 & 初学者向けの説明 \\ \midrule
\endhead
CR & 変更要求。プロジェクト範囲、費用、計画、コード、日程などの変更を求める依頼である。 \\
CCB / SCCB & 構成管理委員会 / ソフトウェア構成管理委員会。変更を承認、修正、延期、却下する権限を持つ。 \\
CI & 構成品目。ひとつの管理単位として扱う品目である。 \\
SCI & ソフトウェア構成品目。ソフトウェアに関する構成品目である。 \\
CM / SCM & 構成管理 / ソフトウェア構成管理。品目を識別し、変更を統制し、状態を記録し、監査する管理活動である。 \\
SCMP & ソフトウェア構成管理計画。構成管理をどう進めるかを定める生きた文書である。 \\
SCSA & ソフトウェア構成状態記録・報告。構成品目、基準版、変更、関係を記録し報告する活動である。 \\
FCA / PCA & 機能構成監査 / 物理構成監査。仕様に合うか、実物と文書が合うかを確認する。 \\
SBOM & ソフトウェア部品表。ビルドに使われた部品と供給網上の関係を記録する。 \\
VDD & 版記述文書。リリースに何が含まれるかを物理的に記述する文書である。 \\
\bottomrule\end{longtable}
\section*{導入：ソフトウェア構成管理が必要な理由}
\addcontentsline{toc}{section}{導入：ソフトウェア構成管理が必要な理由}
ソフトウェア構成管理は、ソフトウェアライフサイクル全体を通じて構成管理を適用し、構成品目の完全性と正確性を保つ活動である。構成管理は、構成品目の機能的・物理的特性を識別して文書化し、それらへの変更を統制し、変更処理と実施状況を記録・報告し、指定要求への適合を検証するための技術的・管理的な規律である。

SCM は、プロジェクト管理、開発、保守、品質保証、顧客、最終利用者を支える。開発方法がウォーターフォール型であっても、反復・増分型であっても、管理対象を明確にし、変更と版を管理し、証跡を残す必要がある点は変わらない。

SCM はソフトウェア品質保証と密接に関係する。品質保証は、成果物とプロセスが要求に合っていることを示す活動である。SCM は、構成品目の識別、変更統制、状態記録、監査を通じて、その証拠を残す。

\noteBox{重要}{SCM の目的は、変更を禁止することではない。変更を安全に行い、何が、誰によって、いつ、なぜ、どの版に対して変わったのかを後から説明できるようにすることである。}
\todobox{SCM が支える活動の関係図を入れる}{白背景、モノクロ、A4 本文幅の概念図を作成する。中央に「SCM」を置き、左に「開発」、右に「保守」、上に「プロジェクト管理」、下に「品質保証」、右下に「顧客・利用者」を配置する。SCM から各活動へ矢印を引き、「識別」「変更統制」「状態記録」「監査」「リリース」の短いラベルを付ける。日本語ラベル、講義ノート風。}
\section*{1 構成管理プロセスの管理}
\addcontentsline{toc}{section}{1 構成管理プロセスの管理}
構成管理プロセスの管理は、製品の進化と完全性を統制するための土台である。SCM は、構成品目を識別し、変更を管理し、構成情報を検証・記録・報告する。成功する SCM には、組織の状況、制約、標準、プロジェクトの性質を理解した計画が必要である。

\subsection*{1.1 組織における構成管理の位置づけ}
\addcontentsline{toc}{subsection}{1.1 組織における構成管理の位置づけ}
SCM は、特定のプロジェクトだけで完結しない。開発組織、保守組織、品質保証組織、運用組織、外部委託先、場合によってはハードウェアやファームウェアの構成管理とも関係する。大規模なシステムでは、ソフトウェア SCM はシステム全体の構成管理と整合していなければならない。

SCM の全体責任は、独立した構成管理部門または指定された個人に置かれることがある。一方で、日々のタスクは開発部門や保守部門に分担されることがある。したがって、誰がどの活動に責任を持つのかを曖昧にしないことが重要である。

\subsection*{1.2 構成管理プロセスに対する制約と指針}
\addcontentsline{toc}{subsection}{1.2 構成管理プロセスに対する制約と指針}
SCM プロセスには、社内方針、手順、契約、規制、標準、採用するライフサイクルモデルが影響する。安全に関わるソフトウェアでは、外部規制機関が構成監査や管理対象品目を要求することがある。

開発方法も SCM に影響する。継続的統合では、頻繁なビルド、テスト、配備が起こるため、分岐、統合、ビルド、試験、成果物保存の手順を明確にする必要がある。

\todobox{分岐・統合と継続的統合の関係図を入れる}{白背景、モノクロ、A4 本文幅の横長図を作成する。左から「開発ブランチ」「変更」「マージ」「ビルド」「自動テスト」「成果物」「リリース候補」へ進む流れを描く。分岐と統合が SCM 計画の対象であることを注記する。日本語ラベル、細線、講義ノート風。}
\subsection*{1.3 構成管理の計画}
\addcontentsline{toc}{subsection}{1.3 構成管理の計画}
構成管理の計画では、構成識別、構成統制、状態記録、構成監査、リリース管理と提供をどのように実施するかを定める。さらに、組織と責任、資源とスケジュール、ツール選定、供給者統制、インタフェース統制を考慮する。

\noteBox{計画時に見る観点}{プロジェクト規模、安全性・セキュリティの重要度、契約・規制、開発拠点の分散、外部委託の有無、使うツール、分岐・統合方針、ビルドとテストの頻度を確認する。}
\subsubsection*{1.3.1 組織と責任}
\addcontentsline{toc}{subsubsection}{1.3.1 組織と責任}
SCM 活動の責任者、承認権限、報告経路を明確にする。役割は個人名だけでなく、部署や役職として定義すると、担当者交代時にも維持しやすい。

\subsubsection*{1.3.2 資源とスケジュール}
\addcontentsline{toc}{subsubsection}{1.3.2 資源とスケジュール}
SCM 活動に必要な人員、ツール、教育、作業順序、節目を計画する。構成管理は後から付け足す活動ではなく、要求、設計、実装、テスト、リリースと並行して計画される。

\subsubsection*{1.3.3 ツール選定と実装}
\addcontentsline{toc}{subsubsection}{1.3.3 ツール選定と実装}
SCM では、単一の万能ツールではなく、版管理、変更管理、ビルド、自動テスト、リポジトリ、リリース管理などの道具群を組み合わせることが多い。ツール選定では、組織方針、費用、保守、移行、将来利用、既存ツールとの統合、分岐・統合戦略への適合を確認する。

ツール群は、異なる供給元のツールを組み合わせる開放型の作業台として構成する場合と、統合製品群として構成する場合がある。小規模組織と大規模組織では、必要な統合範囲が異なる。

\subsubsection*{1.3.4 供給者・外部委託先統制}
\addcontentsline{toc}{subsubsection}{1.3.4 供給者・外部委託先統制}
購入したソフトウェア、コンパイラ、外部ライブラリ、外部委託先が作るソフトウェアも構成管理の対象になり得る。変更や更新をどのように評価し、プロジェクトのライブラリに取り込み、証跡を残すかを計画する。

\subsubsection*{1.3.5 インタフェース統制}
\addcontentsline{toc}{subsubsection}{1.3.5 インタフェース統制}
ソフトウェア品目が他のソフトウェアやハードウェアと接続する場合、一方の変更は他方へ影響する。インタフェース仕様、インタフェース統制計画、インタフェース統制文書を用いて、接続点を識別し、変更を管理し、関係者へ伝える必要がある。

\subsection*{1.4 ソフトウェア構成管理計画}
\addcontentsline{toc}{subsection}{1.4 ソフトウェア構成管理計画}
ソフトウェア構成管理計画（SCMP）は、SCM 計画の結果を記録する生きた文書である。ライフサイクル中に必要に応じて更新・承認される。SCMP だけでは日々の運用に足りないため、分岐方針、ビルド頻度、試験手順などの下位手順を別途定めることが多い。

\begin{longtable}{p{0.30\linewidth}p{0.64\linewidth}}\toprule
SCMP の項目 & 内容 \\ \midrule
\endfirsthead\toprule
SCMP の項目 & 内容 \\ \midrule
\endhead
導入 & 目的、適用範囲、用語を定める。 \\
SCM 管理 & 組織、責任、権限、適用する方針、手順を定める。 \\
SCM 活動 & 構成識別、構成統制、状態記録、監査、リリース管理などを定める。 \\
SCM スケジュール & 他のプロジェクト活動との調整、節目、実施時期を定める。 \\
SCM 資源 & ツール、物理資源、人員、教育を定める。 \\
SCMP 保守 & 計画書をいつ、誰が、どう改訂・承認するかを定める。 \\
\bottomrule\end{longtable}
\subsection*{1.5 ソフトウェア構成管理の監視}
\addcontentsline{toc}{subsection}{1.5 ソフトウェア構成管理の監視}
SCM プロセスを実装した後は、SCMP の規定が正しく実施されているかを監視する。SCM 責任者は、割り当てられたタスクが実施されているかを確認する。品質保証部門が遵守監査の一部として監視する場合もある。

統合された SCM ツールは、手順遵守を支援する。一部のツールは柔軟性を残し、別のツールは特定の手順を強制する。どの程度の柔軟性を許すかは、ツール選定時の重要な判断である。

\subsubsection*{1.5.1 測定}
\addcontentsline{toc}{subsubsection}{1.5.1 測定}
SCM 測定は、製品の進化、プロセスの効率、改善の機会を把握するために使う。変更要求数、変更実施時間、構成品目の状態、監査結果、再作業量などは、構成管理の健全性を示す手がかりになる。

\subsubsection*{1.5.2 プロセス中監査}
\addcontentsline{toc}{subsubsection}{1.5.2 プロセス中監査}
プロセス中監査は、ライフサイクル中に SCM の実施状況を確認する活動である。計画どおりに構成品目が識別されているか、承認済み変更だけが取り込まれているか、状態情報が正しく記録されているかを確認する。

\todobox{SCM 計画・実施・監視の循環図を入れる}{白背景、モノクロ、A4 本文幅の循環図を作成する。「計画」「実施」「測定」「監査」「改善」の五つを円環で配置する。中央に「SCMP」を置き、各活動が SCMP に基づき、結果で SCMP を更新することを示す。日本語ラベル、講義ノート風。}
\section*{2 ソフトウェア構成識別}
\addcontentsline{toc}{section}{2 ソフトウェア構成識別}
ソフトウェア構成識別は、管理すべき品目を識別し、品目と版の識別方式を定め、管理するための道具と技法を決める活動である。構成識別は、変更統制、状態記録、監査、リリース管理の前提になる。

\subsection*{2.1 管理対象品目の識別}
\addcontentsline{toc}{subsection}{2.1 管理対象品目の識別}
変更を統制する最初の手順は、何を管理対象にするかを決めることである。ソフトウェア構成をシステム構成の中で理解し、ソフトウェア構成品目を選び、ラベル付け方針を決める。

\subsubsection*{2.1.1 ソフトウェア構成}
\addcontentsline{toc}{subsubsection}{2.1.1 ソフトウェア構成}
ソフトウェア構成とは、技術文書で定められた、または製品として実現された、ソフトウェアの機能的・物理的特性である。システム全体の構成の一部として理解できる。

\subsubsection*{2.1.2 ソフトウェア構成品目}
\addcontentsline{toc}{subsubsection}{2.1.2 ソフトウェア構成品目}
構成品目は、ひとつの管理単位として扱うハードウェア、ソフトウェア、またはその組合せである。ソフトウェア構成品目は、構成品目として確立されたソフトウェア実体である。

管理対象になり得る品目には、計画書、仕様書、設計文書、テスト資料、ソフトウェアツール、ソースコード、実行コード、コードライブラリ、データ、データ辞書、インストール・保守・運用・利用者向け文書が含まれる。

\noteBox{選定のバランス}{品目を細かくしすぎると管理量が増え、粗くしすぎると変更の影響が見えにくい。適切な可視性と管理可能な数の間でバランスを取る。}
\todobox{構成品目の範囲を示す図を入れる}{白背景、モノクロ、A4 本文幅の集合図を作成する。大きな枠を「ソフトウェア構成」とし、内側に「要求・仕様」「設計文書」「ソースコード」「テスト資料」「データ」「外部ライブラリ」「ビルドツール」「利用者文書」を配置する。構成品目として管理対象に選ぶ品目にチェック印を付ける。日本語ラベル、講義ノート風。}
\subsection*{2.2 構成品目の識別子と属性}
\addcontentsline{toc}{subsection}{2.2 構成品目の識別子と属性}
構成品目には、一意に識別できる名前や番号を付ける。状態記録・報告は、これらの属性を集めることで実現される。

\begin{longtable}{p{0.30\linewidth}p{0.64\linewidth}}\toprule
属性 & 説明 \\ \midrule
\endfirsthead\toprule
属性 & 説明 \\ \midrule
\endhead
CI 名 & 人が読んで意味を理解しやすい名称である。 \\
一意識別子 & システム上で品目を一意に指す記号や番号である。 \\
説明 & 品目の目的、内容、利用範囲を説明する。 \\
日付 & 作成日、変更日、承認日などを記録する。 \\
種類 & 文書、コード、テスト、データ、外部ライブラリなどの分類である。 \\
所有者 & 品目の管理責任を持つ人または組織である。 \\
版 & どの時点の状態かを示す識別情報である。 \\
\bottomrule\end{longtable}
識別子には、意味を持たせる方法と、単なる連番のように意味を持たせない方法がある。意味を持たせる場合、反復番号、品目分類、連番などを組み合わせることがある。ただし、意味を持たせすぎると、分類変更時に識別子が扱いにくくなる。

\subsection*{2.3 基準版の識別}
\addcontentsline{toc}{subsection}{2.3 基準版の識別}
基準版とは、ある時点で正式に承認され、固定された構成品目またはその集合である。基準版は、正式な変更統制手続きなしには変更できない。承認済み変更を含む基準版は、その時点の承認済み構成を表す。

基準版は、開発を止めるためのものではない。むしろ、変更の基点を明確にするためのものである。「この要求書の版に基づいて設計した」「このソースコードの版からリリースした」「この基準版に対して不具合修正を入れた」と説明できるようにする。

\subsection*{2.4 基準版の属性}
\addcontentsline{toc}{subsection}{2.4 基準版の属性}
基準版にも属性を持たせる。代表的な属性は、基準版名、一意識別子、説明、作成日、含まれる構成品目である。これにより、後から「どの基準版が何を含んでいたか」を確認できる。

\todobox{基準版と変更統制の関係図を入れる}{白背景、モノクロ、A4 本文幅の横長図を作成する。左に「作業中の品目」、中央に「承認」、右に「基準版」を置く。基準版から下方向に「変更要求」「影響分析」「承認」「変更実施」「新しい基準版」へ進む矢印を描く。基準版は正式手続きなしに変更できないことを注記する。日本語ラベル、講義ノート風。}
\subsection*{2.5 関係方式の定義}
\addcontentsline{toc}{subsection}{2.5 関係方式の定義}
構成品目は単独ではなく、他の品目と関係を持つ。関係を記録すると、変更要求の影響分析、版の追跡、部品表の作成が容易になる。代表的な関係は、依存、派生、継承、変種である。

\begin{longtable}{p{0.30\linewidth}p{0.64\linewidth}}\toprule
関係 & 説明 \\ \midrule
\endfirsthead\toprule
関係 & 説明 \\ \midrule
\endhead
依存 & 二つの構成品目が互いに影響し合う関係である。例として、クラスモデルとシーケンス図の整合がある。 \\
派生 & ある構成品目から別の構成品目が生じる関係である。例として、仕様から設計、設計からコードが派生する。 \\
継承 & 同じ品目が時間とともに版として進化する関係である。版 1 から版 2 への移り変わりを表す。 \\
変種 & 意図的に用意された代替版である。プラットフォーム別、機能別、顧客別の版が例である。 \\
\bottomrule\end{longtable}
どの関係を追跡するかは慎重に決める。追跡すれば変更判断に役立つが、記録と保守の作業が増える。ソフトウェア部品表は、ビルドに使われる構成品目とその供給網上の関係を記録する正式な手段である。

\todobox{構成品目間の関係図を入れる}{白背景、モノクロ、A4 本文幅のネットワーク図を作成する。ノードとして「要求仕様」「設計モデル」「ソースコード」「外部ライブラリ」「テストケース」「実行ファイル」を置く。矢印に「依存」「派生」「継承」「変種」のラベルを付ける。右下に「SBOM は部品と関係を記録する」と注記する。日本語ラベル、講義ノート風。}
\subsection*{2.6 ソフトウェアライブラリ}
\addcontentsline{toc}{subsection}{2.6 ソフトウェアライブラリ}
ソフトウェアライブラリは、ソースコード、スクリプト、オブジェクトコード、文書、関連成果物を統制された形で保存する集合である。要求やテストケースはリポジトリに保存し、コードの基準版と関連付けるべきである。

ソースコードは版管理システムに保存され、追跡可能性とセキュリティを支える。ビルドで得られるバイナリはリポジトリに保存され、ハッシュなどにより物理構成監査を支援する。決定版媒体ライブラリには、試験、ステージング、本番へ配備できる成果物のリリース基準版が含まれる。

ライブラリ管理では、アクセス制御とバックアップが重要である。リリース管理は、配備される成果物をこれらのライブラリから正しく選ぶことに依存する。

\section*{3 ソフトウェア構成変更統制}
\addcontentsline{toc}{section}{3 ソフトウェア構成変更統制}
ソフトウェア構成変更統制は、ライフサイクル中に構成品目へ必要となる変更を扱う活動である。何を変更するか、誰が承認するか、どのように実施するか、要求からの逸脱や免除をどう扱うかを定める。

変更統制から得られる情報は、変更量、破壊的変更、再作業、欠陥修正、拡張要求の測定にも使える。

\subsection*{3.1 ソフトウェア変更の要求・評価・承認}
\addcontentsline{toc}{subsection}{3.1 ソフトウェア変更の要求・評価・承認}
管理対象品目への変更は、通常、変更要求として提出される。ソフトウェア変更要求プロセスは、変更要求の提出、記録、費用と影響の評価、承認・修正・延期・却下を扱う正式な手順である。

変更要求は、ライフサイクルのどの時点でも、誰からでも発生し得る。問題報告に対する是正処置、機能追加、範囲変更、スケジュール変更、費用変更、手順変更などが原因になり得る。

変更要求を受け取ると、技術評価または影響分析を行う。どの構成品目が影響を受けるか、どの設計・コード・テスト・文書を変更する必要があるかを調べる。ここで、構成品目間の関係情報が重要になる。

\todobox{変更要求プロセスの流れを入れる}{白背景、モノクロ、A4 本文幅のフローチャートを作成する。「変更の必要」「SCR 作成」「予備調査」「CCB レビュー」「承認 / 修正 / 延期 / 却下」「担当割当」「設計・実装・テスト」「変更完了」「通知」を順に配置する。緊急時には実施後に正式手続きを行う例外経路を点線で示す。日本語ラベル、講義ノート風。}
\subsubsection*{3.1.1 ソフトウェア構成管理委員会}
\addcontentsline{toc}{subsubsection}{3.1.1 ソフトウェア構成管理委員会}
変更の承認または却下の権限は、構成管理委員会に置かれる。小規模プロジェクトでは、委員会ではなくプロジェクトリーダーや指定された個人が権限を持つこともある。変更の重要性、予算・日程への影響、ライフサイクル上の時点によって、承認権限の階層を分ける場合がある。

委員会の範囲がソフトウェアに限られる場合は、ソフトウェア構成管理委員会と呼ぶ。SCM 代表者は常に関与し、必要な利害関係者が委員会レベルに応じて参加する。

\subsubsection*{3.1.2 ソフトウェア変更要求プロセス}
\addcontentsline{toc}{subsubsection}{3.1.2 ソフトウェア変更要求プロセス}
有効な SCR プロセスには、変更要求の発行、ワークフローの強制、委員会判断の記録、変更処理情報の報告を支援する手順とツールが必要である。問題報告システムと連携すると、問題解決の追跡と解決速度の把握がしやすい。

\subsubsection*{3.1.3 ソフトウェア変更要求フォーム}
\addcontentsline{toc}{subsubsection}{3.1.3 ソフトウェア変更要求フォーム}
変更要求には、要求の内容と根拠を示すフォームが必要である。承認された場合には、変更の認証や完了確認に使うフォームも必要になる。ツールで管理する場合でも、何を入力し、誰が責任を持つかは人が設計する。

\subsection*{3.2 ソフトウェア変更の実施}
\addcontentsline{toc}{subsection}{3.2 ソフトウェア変更の実施}
承認された変更要求は、定義済みの手順とスケジュールに従って実施する。複数の変更が同時に実施される場合、どの変更要求がどの版や基準版に含まれたかを追跡できなければならない。

変更完了時には、構成監査や品質検証を行い、承認済み変更だけが行われたことを確認する。版管理ツールは、チェックアウト、変更、記録、保存、差し戻しを支援する。並行開発や分散開発では、より強力なツールが必要になる。

\subsection*{3.3 逸脱と免除}
\addcontentsline{toc}{subsection}{3.3 逸脱と免除}
逸脱は、品目を作る前に、特定の性能要求または設計要求から一時的に外れることを認める書面上の許可である。免除は、生産中または検査後に見つかった問題に対して、指定要求から外れていても使用に適すると判断して受け入れる書面上の許可である。

\begin{longtable}{p{0.20\linewidth}p{0.25\linewidth}p{0.49\linewidth}}\toprule
区分 & 時点 & 意味 \\ \midrule
\endfirsthead\toprule
区分 & 時点 & 意味 \\ \midrule
\endhead
逸脱 & 品目を作る前 & 要求を満たせないことが事前にわかり、一定範囲で外れることを認める。 \\
免除 & 生産中または検査後 & 要求から外れた品目を、そのまままたは承認済み手直し後に使用可とする。 \\
\bottomrule\end{longtable}
\section*{4 ソフトウェア構成状態記録・報告}
\addcontentsline{toc}{section}{4 ソフトウェア構成状態記録・報告}
ソフトウェア構成状態記録・報告は、構成品目、基準版、構成品目間の関係について、構成を有効に管理するために必要な情報を記録し、報告する活動である。これは、構成識別で定めた論理的方式に従って実施する。

\subsection*{4.1 ソフトウェア構成状態情報}
\addcontentsline{toc}{subsection}{4.1 ソフトウェア構成状態情報}
SCSA は、ライフサイクルの進行に伴って必要な情報を取得し、検証し、妥当性を確認し、報告する仕組みを設計・運用する。必要に応じて、状態情報そのものも保護する必要がある。

\begin{itemize}
\item 現在有効な構成識別と承認済み構成識別。
\item 変更の現在の実施状況。
\item 影響を受ける構成品目と関連システム。
\item 逸脱と免除。
\item 検証・妥当性確認活動の証跡。
\item 完全性の指標、セキュリティ状態、基準版状態、構成品目ごとの変更要求数、変更実施の平均時間。
\end{itemize}
\subsection*{4.2 ソフトウェア構成状態報告}
\addcontentsline{toc}{subsection}{4.2 ソフトウェア構成状態報告}
報告情報は、開発チーム、運用、セキュリティ、保守、プロジェクト管理、品質保証など、さまざまな組織要素に使われる。報告には、自動レポート、特定質問に対する随時照会、定期レポートがある。

SCSA がライフサイクル全体で生む情報は、品質保証、セキュリティ、規制、ガバナンスへの適合を示す証拠になる。したがって、状態記録は単なる管理表ではなく、説明責任を支える証跡である。

\todobox{状態記録・報告が証跡になる流れを入れる}{白背景、モノクロ、A4 本文幅の横長図を作成する。左から「構成品目」「変更」「状態記録」「報告」「品質保証・セキュリティ・規制対応」へ進む流れを描く。状態記録の下に「誰が」「いつ」「何を」「どの版へ」「どの承認で」を並べる。日本語ラベル、講義ノート風。}
\section*{5 ソフトウェア構成監査}
\addcontentsline{toc}{section}{5 ソフトウェア構成監査}
ソフトウェア監査は、成果物または成果物集合を独立に調べ、仕様、標準、契約、法規制などの基準に適合しているかを評価する活動である。構成監査は、構成品目が機能的・物理的特性の要求をどの程度満たしているかを確認する。

構成監査は、ライフサイクル中の重要な時点で非公式に行われることもある。契約や安全上の要求により、正式な機能構成監査と物理構成監査が必要になることもある。

\subsection*{5.1 ソフトウェア機能構成監査}
\addcontentsline{toc}{subsection}{5.1 ソフトウェア機能構成監査}
ソフトウェア機能構成監査は、監査対象のソフトウェア品目が、支配的な仕様と整合していることを確認する。検証・妥当性確認活動の出力は、この監査の重要な入力になる。

\subsection*{5.2 ソフトウェア物理構成監査}
\addcontentsline{toc}{subsection}{5.2 ソフトウェア物理構成監査}
ソフトウェア物理構成監査は、設計文書や参照文書が、実際に作られたソフトウェア製品と整合していることを確認する。簡単に言えば、機能構成監査は「仕様どおりに機能するか」、物理構成監査は「実物と文書が一致するか」を見る。

\begin{longtable}{p{0.20\linewidth}p{0.25\linewidth}p{0.49\linewidth}}\toprule
監査 & 主に見ること & 代表的な入力 \\ \midrule
\endfirsthead\toprule
監査 & 主に見ること & 代表的な入力 \\ \midrule
\endhead
機能構成監査 & 仕様どおりの機能・性能を満たすか。 & 要求、仕様、検証・妥当性確認結果。 \\
物理構成監査 & 実物、設計文書、参照文書、リリース内容が一致するか。 & 設計文書、ソース、バイナリ、ハッシュ、版記述文書。 \\
\bottomrule\end{longtable}
\subsection*{5.3 ソフトウェア基準版のプロセス中監査}
\addcontentsline{toc}{subsection}{5.3 ソフトウェア基準版のプロセス中監査}
プロセス中監査は、開発途中で特定の構成要素の状態を調べるために行う。すべての基準版品目について、性能が仕様と整合しているか、発展中の文書が基準版品目と整合しているかを確認できる。

構成識別で見つけた構成品目を継続的にレビューすると、ガバナンスや規制要求への適合を確認しやすい。構成監査レビューは、構成品目をレビューすべき時点ごとに、開発全体を通じて行われる。

\todobox{機能構成監査と物理構成監査の違いを示す図を入れる}{白背景、モノクロ、A4 本文幅の二分割図を作成する。左に「機能構成監査：仕様と振る舞いの一致」、右に「物理構成監査：実物と文書の一致」を置く。左には「要求」「仕様」「V\&V 結果」から監査へ向かう矢印、右には「設計文書」「ソース」「バイナリ」「VDD」「ハッシュ」から監査へ向かう矢印を描く。日本語ラベル、講義ノート風。}
\section*{6 ソフトウェアリリース管理と提供}
\addcontentsline{toc}{section}{6 ソフトウェアリリース管理と提供}
ここでいうリリースは、開発活動の外へソフトウェアと関連成果物を配布することである。社内リリースや顧客への配布を含む。複数のプラットフォーム版や機能差のある版がある場合、特定版の再現、正しい材料の包装、正しい対象への提供が必要になる。

\subsection*{6.1 ソフトウェアビルド}
\addcontentsline{toc}{subsection}{6.1 ソフトウェアビルド}
ソフトウェアビルドは、正しい版のソフトウェア構成品目と適切な構成データを使って、顧客やテストチームなどに提供するソフトウェアパッケージを作る活動である。ビルド手順は、正しい順序で手順を実施するために必要である。

以前のリリースを正確に再構築できることも重要である。そのためには、ソースコードだけでなく、コンパイラ、ビルドツール、スクリプト、設定、ビルド手順も SCM の管理下に置く必要がある。

継続的統合では、変更がソース管理リポジトリへコミットされると、ビルドが自動的に実行される。パイプラインは、取得、コンパイル、リンク、静的解析、単体テスト、コード網羅率測定、文書抽出などを行う。ビルドに含まれる成果物を示す SBOM は重要な構成管理出力である。

\subsection*{6.2 ソフトウェアリリース管理}
\addcontentsline{toc}{subsection}{6.2 ソフトウェアリリース管理}
ソフトウェアリリース管理は、実行プログラム、文書、リリースノート、構成データなど、製品要素を識別し、包装し、提供する活動である。いつリリースするかは、解決した問題の重大度、過去リリースの欠陥密度、顧客通知の必要性などによって判断する。

リリースの物理的内容を記録する情報は、版記述文書である。リリースノートには、新機能、既知の問題、必要なプラットフォーム条件を記す。パッケージには、インストールまたはアップグレード手順も含める。完全性を保証するために、デジタル署名などを使うこともある。

継続的デリバリでは、ビルド成果物をインストールパッケージへまとめ、検証環境または本番環境へ配備するためのパイプラインを設ける。

\todobox{ビルドからリリースまでのパイプライン図を入れる}{白背景、モノクロ、A4 本文幅の横長パイプライン図を作成する。左から「ソース管理」「ビルド」「静的解析」「単体テスト」「成果物リポジトリ」「SBOM」「VDD」「リリースノート」「署名」「配備」へ進む。下段に「再現可能性」「証跡」「完全性」の注記を置く。日本語ラベル、講義ノート風。}
\section*{7 ソフトウェア構成管理ツール}
\addcontentsline{toc}{section}{7 ソフトウェア構成管理ツール}
構成管理ツールは、構成管理活動を支援する技術と論理を提供する。SCM は他のプロセスや既存ツールと統合されるほど効果が高まる。ツールの範囲は、利用者、組織規模、製品の変種、規制要求、開発分散の程度によって変わる。

\begin{longtable}{p{0.30\linewidth}p{0.64\linewidth}}\toprule
ツール分類 & 主な役割 \\ \midrule
\endfirsthead\toprule
ツール分類 & 主な役割 \\ \midrule
\endhead
版管理ツール & ソースコード、設定、文書などの版を保存し、変更履歴を管理する。 \\
ビルド自動化ツール & コンパイル、リンク、品質確認、テスト、成果物生成を自動化する。 \\
変更統制ツール & 変更要求の登録、ワークフロー、承認、通知、状態追跡を支援する。 \\
リポジトリ管理 & ビルドで生成されたバイナリ、ライブラリ、パッケージを保存する。 \\
CMDB / 永続化基盤 & 構成情報、関係、状態、証跡を保存する。 \\
リリース・配備ツール & 成果物の包装、配布、配備、ロールバック、署名を支援する。 \\
\bottomrule\end{longtable}
小規模な組織や、製品の変種を出さない開発グループでは、個別支援ツールで十分なことがある。大規模組織や複雑な製品では、ワークフロー、役割、承認要求まで含む全社的な支援ツールが必要になる。

\todobox{構成管理ツールの範囲を示すピラミッドを入れる}{白背景、モノクロ、A4 本文幅の縦長ピラミッド図を作成する。下段に「個別支援ツール：版管理、ビルド、変更管理」、中段に「プロジェクト支援ツール：作業空間、分散開発、並行開発」、上段に「全社プロセス支援ツール：ワークフロー、役割、認証要求」を配置する。右側に「範囲と形式性が上がる」と矢印を描く。日本語ラベル、講義ノート風。}
\section*{8 トピックと参考文献の対応}
\addcontentsline{toc}{section}{8 トピックと参考文献の対応}
この表は、SWEBOK が示す「トピック対参考資料の行列」を、学習用に簡略化したものである。詳しく学ぶときは、各トピックに対応する文献を読む。

\begin{longtable}{p{0.30\linewidth}p{0.64\linewidth}}\toprule
トピック & 主な参考資料 \\ \midrule
\endfirsthead\toprule
トピック & 主な参考資料 \\ \midrule
\endhead
1. 構成管理プロセスの管理 & IEEE Std 828; Hass; Sommerville \\
1.1 組織における位置づけ & IEEE Std 828; Hass; Sommerville \\
1.2 制約と指針 & IEEE Std 828; Hass; Moore \\
1.3 構成管理の計画 & IEEE Std 828; Hass; Sommerville \\
1.3.3 ツール選定と実装 & Hass \\
1.4 ソフトウェア構成管理計画 & IEEE Std 828; Hass \\
1.5 構成管理の監視 & Hass \\
2. ソフトウェア構成識別 & IEEE Std 828; Hass \\
2.1 管理対象品目の識別 & IEEE Std 828; Hass \\
2.1.1 ソフトウェア構成 & ISO/IEC/IEEE 24765; IEEE Std 828 \\
2.1.2 ソフトウェア構成品目 & ISO/IEC/IEEE 24765; Hass \\
2.2 識別子と属性 & IEEE Std 828; Hass \\
2.5 関係方式の定義 & Hass \\
2.6 ソフトウェアライブラリ & IEEE Std 828; Hass \\
3. ソフトウェア構成変更統制 & IEEE Std 828; Hass; Sommerville; Moore \\
3.1 変更の要求・評価・承認 & Hass; Sommerville; Moore \\
3.2 ソフトウェア変更の実施 & Sommerville; Moore \\
4. 構成状態記録・報告 & IEEE Std 828; Hass; Moore \\
5. ソフトウェア構成監査 & IEEE Std 828; Moore \\
6. リリース管理と提供 & IEEE Std 828; Hass; Sommerville \\
7. 構成管理ツール & Hass \\
\bottomrule\end{longtable}
\section*{9 発展的な読み物}
\addcontentsline{toc}{section}{9 発展的な読み物}
SWEBOK 第8章で紹介される発展的な読み物を、日本語で読む目的に合わせて整理する。

\begin{longtable}{p{0.30\linewidth}p{0.64\linewidth}}\toprule
文献 & 読むと得られること \\ \midrule
\endfirsthead\toprule
文献 & 読むと得られること \\ \midrule
\endhead
Berczuk and Appleton, Software Configuration Management Patterns & 構成管理の実践を「パターン」として学べる。特定ツールに依存せず、チーム作業、統合、分岐、変更管理を考える助けになる。 \\
CMMI for Development & 成熟度レベル 2 の構成管理活動を含む。組織のプロセス改善と構成管理を結びつけて理解できる。 \\
Aiello and Sachs, Configuration Management Best Practices & 現実の組織で使える構成管理の実践を扱う。変更統制の種類や運用上の注意点を深く学べる。 \\
\bottomrule\end{longtable}
\section*{10 参考文献}
\addcontentsline{toc}{section}{10 参考文献}
1. ISO/IEC/IEEE, ISO/IEC/IEEE 24765:2017 Systems and Software Engineering — Vocabulary, 2nd ed., 2017.

2. IEEE, IEEE Standard 828-2012, Standard for Configuration Management in Systems and Software Engineering, 2012.

3. A. M. J. Hass, Configuration Management Principles and Practices, 1st ed., Addison-Wesley, 2003.

4. I. Sommerville, Software Engineering, 10th ed., Pearson, 2016.

5. J. W. Moore, The Road Map to Software Engineering: A Standards-Based Guide, Wiley-IEEE Computer Society Press, 2006.

6. S. P. Berczuk and B. Appleton, Software Configuration Management Patterns: Effective Teamwork, Practical Integration, Addison-Wesley, 2003.

7. CMMI Institute, CMMI for Development, Version 2.0, 2018.

8. B. Aiello and L. A. Sachs, Configuration Management Best Practices: Practical Methods That Work in the Real World, 2011.

\section*{11 画像 TODO 一覧}
\addcontentsline{toc}{section}{11 画像 TODO 一覧}
本文に配置した画像 TODO を再掲する。実際に図を作る場合は、各 TODO のプロンプトをそのまま画像生成に使うと、A4 資料に合う図を作りやすい。

\begin{itemize}
\item 第8章の全体地図。
\item SCM が支える活動の関係図。
\item 分岐・統合と継続的統合の関係図。
\item SCM 計画・実施・監視の循環図。
\item 構成品目の範囲を示す図。
\item 基準版と変更統制の関係図。
\item 構成品目間の関係図。
\item 変更要求プロセスの流れ。
\item 状態記録・報告が証跡になる流れ。
\item 機能構成監査と物理構成監査の違いを示す図。
\item ビルドからリリースまでのパイプライン図。
\item 構成管理ツールの範囲を示すピラミッド。
\end{itemize}
\section*{12 章末確認問題}
\addcontentsline{toc}{section}{12 章末確認問題}
\begin{enumerate}
\item ソフトウェア構成管理が、単なる版管理ツールの利用ではない理由を説明せよ。
\item 構成品目とソフトウェア構成品目の違いを説明せよ。
\item 基準版とは何か。なぜ正式な変更統制なしに変更してはいけないのか。
\item 変更要求、ソフトウェア変更要求、構成管理委員会の関係を説明せよ。
\item 逸脱と免除の違いを説明せよ。
\item 構成状態記録・報告が、品質保証、セキュリティ、規制対応に役立つ理由を説明せよ。
\item 機能構成監査と物理構成監査の違いを説明せよ。
\item ソフトウェアビルドを再現可能にするために、どのような品目を構成管理下に置くべきか。
\item ソフトウェア部品表が、変更影響分析やセキュリティ対応に役立つ理由を説明せよ。
\item 小規模プロジェクトと全社的な開発組織では、構成管理ツールに求める範囲がどう変わるか。
\end{enumerate}
\noteBox{解答の方向性}{1 は「識別、変更統制、状態記録、監査、リリース管理を含む」ことを述べる。2 は「一つの管理単位」と「ソフトウェア関連の管理単位」の関係で説明する。3 は「承認済み構成を示す基点」であることを述べる。6 は「証跡を残す活動」であること、7 は「仕様どおりの機能か」と「実物と文書が一致するか」の違いで説明できる。}

\end{document}
